% ============================================================
% Part III — Desktop Section — M2 (The World and Its Deployment) — DRAFT v1
% Mode: Stage 3 guided drafting (per-paragraph protocol). Plan approved
% 2026-07-20: Batch A = M2.1 artifact (3 ¶¶); Batch B = M2.2 cohort/instrument
% (2 ¶¶); M2.3 method-of-reading BLOCKED by Open decision #1 (arc placement).
% Conventions: as DESKTOP-M0-DRAFT-v2.tex (MLA; abbreviated-title parentheticals;
%   first person for K&S work; register + bans per memory).
% NEW SOURCE: King & Salvo, "Scalable Virtual Reality Global Clinical Immersion
%   for Culturally Responsive Engineering Design Skills" (ASEE paper #49715;
%   venue/year VERIFICATION PENDING via index agent; PDF pages 1-15, no journal
%   pagination — parenthetical cites use PDF pages). Platform facts pp. 3-4;
%   results p. 5; discussion p. 9. Ingestion + vle-05 unit creation running in
%   background agent.
% AUTHOR-CONFIRM FLAGS (Batch A): (1) RESOLVED (AUTHOR-SUPPLIED 2026-07-29,
%   assembly pass): deployed-build procedure list = Bilateral Scar Revision ·
%   PeriMortem C-Section Simulation · Transcatheter Aortic Valve Replacement ·
%   Deep Brain Stimulation · Spinal Deformation Surgery · Pediatric Auricular
%   Reconstruction · Cornea Transplant · Standard Cataract Procedures
%   (Paraguay) · Standard Cataract Procedures (Vietnam) · Standard and Laser
%   Cataract Procedures (UCI); each procedure includes its supplementary
%   video materials (animated models of procedure and/or equipment;
%   interviews with physicians and medical staff). Integrated into ¶1
%   2026-07-29 — REVIEW RULINGS (user 2026-07-31, batch 3): (i) CORRECTED —
%   UCI MC + CHOC cover SEVEN domestic items; the perimortem C-section
%   SIMULATION was recorded at the UCI Medical Education Simulation Center
%   (medschool.uci.edu, Division of Educational Technology; name verified
%   from the center's page). ¶1 restructured accordingly. (ii) CONFIRMED:
%   keep "spinal deformation surgery";
%   (2) RESOLVED (AUTHOR-CONFIRMED 2026-07-24): the launching-point/portal
%   structure was PILOT-ERA ONLY (rudimentary hospital). The deployed build
%   used by the 241 students = the full multi-story hospital; each procedure
%   room contains its own media player; the student walks in and plays the
%   video in place — no teleporting. Capture-verified 2026-07-24 against the
%   three videos at D:\PhD\Dissertation\Media\VLE (frames + transcripts in
%   corpus/index/Wendt.../bridge-sources/bme-vr-env-analysis/);
%   (3) RESOLVED (AUTHOR-CONFIRMED 2026-07-24): the deployed content library
%   STILL CONTAINS the animation and overview items (pilot-paper TAVR
%   inventory: "an animation of a TAVR procedure" + Edwards SAPIEN XT valve
%   overview; SAPIEN XT overview visible in-world in the 2024 group capture)
%   — outline M2.1's "device animations" stands in the artifact inventory;
%   (4) RESOLVED (AUTHOR-CONFIRMED 2026-07-20): Vietnam/Paraguay content was
%   present in ALL THREE surveyed terms.
% M2.1 TODO (AUTHOR-DIRECTED 2026-07-24): add a description of the video
%   delivery system — footage streamed from AWS S3 through the in-world
%   player, adaptive bitrate, ~3-second buffer on play and after scrubbing;
%   player opened with E, content under Procedures/Interviews tabs; 360°
%   footage with physician-POV picture-in-picture inset + Hide-POV toggle
%   (source: MediaPlayer.mp4 how-to narration + group-session transcript;
%   YouTube-streaming rejection reason CORRECTED, AUTHOR-CONFIRMED 2026-07-24:
%   it was NOT POSSIBLE to include YouTube streaming inside the UE
%   environment — a technical impossibility, not a weighed preference).
%   Related (AUTHOR-CONFIRMED 2026-07-24): YouTube-style conveniences such as
%   increased playback speed were DELIBERATELY omitted as a pedagogical
%   decision — feeds M3's prediction context.
% M2.1 (Batch A) — plan approved 2026-07-20; DRAFTED TEXT IN REVIEW (user
% reviews A + B in one pass).
% ============================================================

\section{The World and Its Deployment}

% --- M2.1 The artifact (Batch A, 2026-07-20 — in review)

The deployed platform is a downloadable application for Windows PCs, built in
Unreal Engine, that renders a navigable virtual hospital.\footnote{The platform
facts in this part are reported in a second of my publications, co-authored
with Prof.\ King and Kadin Diec, and cited parenthetically by abbreviated title
hereafter: King, Christine E., Kadin Diec, and Dalton Salvo. ``Scalable Virtual
Reality Global Clinical Immersion for Culturally Responsive Engineering Design
Skills.'' 2025 ASEE Annual Conference \& Exposition, paper \#49715.
% CITATION STATUS: year 2025 VERIFIED (title-page ASEE copyright line);
% conference name INFERRED from Paper-ID + copyright (not printed in body).
% Kadin Diec CONFIRMED as co-author by user 2026-07-20, spelling correct.
The paper has no journal pagination; parenthetical page numbers refer to the
PDF.} Students move through it in the first person, by mouse and keyboard, on
their own monitors. Each procedure room gathers the materials for one
operation or category of operations. The room's principal material is the
operating-room footage, which was captured two ways at once: a 360-degree
stereoscopic camera records the whole room, and a head-mounted camera worn by
the operating physician records the surgeon's own first-person view
(``Scalable Virtual Reality'' 4). Each room also holds recorded interviews
with the physicians and engineers involved and supplementary lectures and
readings (4), together with animated models of the procedure or its
equipment. The filmed procedures now span multiple institutions and
countries. Seven procedures were recorded at the UCI Medical Center and the Children's
Hospital of Orange County: bilateral scar revision, transcatheter aortic
valve replacement, deep brain stimulation, spinal deformation surgery,
pediatric auricular reconstruction, cornea transplant, and standard and laser
cataract procedures. A perimortem C-section simulation was recorded at the
UCI Medical Education Simulation Center. To these the platform adds standard
cataract procedures filmed in hospitals in Vietnam and Paraguay, placing the
same operation side by side in three healthcare systems (3--4). On an ordinary monitor the footage plays as an interactive 360-degree
panorama; the same materials are mirrored to a YouTube channel and website for
students whose machines cannot run the application (4).

The world is built for company. Voice chat is real-time and spatial---students
hear the classmates near them in the virtual space---and a group can convene
inside the environment, watch a procedure together, and share various other
artifacts through the screen sharing functionality: the design ``supported
real-time voice chat and screen sharing, enabling collaborative team viewing if
students wanted to meet in groups within the virtual environment'' (``Scalable
Virtual Reality'' 4). That design choice answers the firmest requirement the
collaborative-learning literature states: in van der Meer and colleagues'
review of 139 virtual-reality collaborative-learning systems, ``efficient
communicational tools are essential for effective collaboration: verbal
(audio) communication is crucial'' (9). The voice channel, that is, is not
one feature among others; it is the medium of speaking-with-others, the mode
in which a group is a group at all.\footnote{Heidegger, reading Aristotle's
definition of the human being as the living thing that speaks:
being-with-one-another ``is not being-with-one-another in the sense of
being-situated-alongside-one-another, but rather in the sense of
\textit{being-as-speaking-with-one-another} through communicating, refuting,
confronting'' (\textit{Basic Concepts} 33; italics Heidegger's).} Every
student is assigned to a procedure-paired group at the start of the term, the
capability is announced in the course materials, and meeting in the world is
encouraged. Group interaction inside the platform is one of the course's stated
practice objectives (3). Co-watching in the virtual hospital is thus not an
incidental feature but the designed use of the environment's social
architecture.

Every element of this environment is authored. The hospital and its rooms, the
two camera vantages, the interviews, the group assignments, and the channel
over which students speak are design decisions, and each is a made source---an
\textit{arch\=e}---of the experience it occasions. The virtual hospital is, in
the dissertation's terms, a manufactured and curated rhetorical ecology: a
designed \textit{Umwelt} whose makers supply the matter and the form of the
encounter while the act remains the student's own. The inventory above is
therefore also an inventory of design levers, and the analysis will return to
each element as the survey data show what it did and did not occasion.

% --- M2.2 The cohort and the instrument (Batch B, 2026-07-20 — in review)
% Facts verified against survey-analysis/00-data-audit, 01-quant-summary,
% 02-codebook-final: N per term 102/79/60; 14 items, 4 open (Q5/Q6/Q9/Q10);
% 873 open responses; strata APP 52/51/50% vs VID 45/46/43% per term,
% chi-sq=0.02, Cramér's V=0.009 (n=231 excl. small UNK) -> term-stable;
% reliability = batch 6, 143 responses double-coded, mean Jaccard .878,
% 4 manual adjudications; codebook v1 unrevised through coding.
% METHODS-TRANSPARENCY: RESOLVED 2026-07-23 (coding review Stations 6+12) —
% disclosure ruling (d) hybrid: plain in-prose clause + substantive footnote at
% the reliability sentence. ¶2 below FINALIZED per Station 12 approval.
% Post-review facts (2026-07-23): ALL canonical passes Fable 5 (Haiku batches
% 1/2/5 re-passed + replaced; cross-model Jaccard .596/.514/.668 documented);
% 7 manual adjudications; α-MASI .82 companion; per-code agreement 50–100%.
% Records: survey-analysis/02, 07-repass-impact-memo, coding/xmodel-*.
% Minimal coding-methods appendix APPROVED — draft at
% drafts/DESKTOP-APPENDIX-CODING-METHODS-DRAFT-v1.tex.

The survey corpus comes from three offerings of the course: Winter 2025 (102
respondents), Spring 2025 (79), and Winter 2026 (60)---241 students in all.
Each term used the same 14-item end-of-course instrument, four items of which
were open-ended; those four items produced the 873 open-text responses analyzed
in this section. One wording note applies throughout: the instrument asks about
the ``VR platform,'' phrasing inherited from the pilot era, though the
deployment it surveys is the 2D-PC-primary build described above. The
instrument also recorded whether students could run the application and whether
the platform elicited immersion, and those answers divide each cohort into two
strata: students who used the desktop application (52\%, 51\%, and 50\% across
the three terms) and students who watched the videos only (45\%, 46\%, and
43\%). The split is statistically indistinguishable across terms
($\chi^2 = 0.02$, Cram\'er's $V = 0.009$, $n = 231$), and that stability
carries weight in what follows: it licenses reading differences between the
strata as structural features of the deployment rather than accidents of a
given cohort. The course survey is itself anchored in the global deployment:
its title is ``Global Healthcare Needs and Final Assessment Survey,'' and its
three content items test the global material---nearly all students answered
them correctly. The dedicated global-procedures questionnaire---two Likert
items on exploring healthcare and on seeing cross-country differences, and an
open invitation to comment on the Vietnam and Paraguay videos---was a separate
Winter 2025 administration, reported in the ASEE paper above; two of its items
share their wording with this survey's platform-use and immersion questions.
Where its results are used, they are cited to that paper.
% [M2.2 ¶1 global-questions revision APPROVED 2026-07-20 — verified against
%  raw Canvas headers (all three terms) + ASEE paper Appendix B.]

% ****** USER MANUAL-REWRITE PENDING (ruled 2026-07-31, readthrough batch 3):
% ****** the author will personally rewrite THIS PARAGRAPH ("Responses were
% ****** anonymized...") AND its coding-passes footnote ("All coding
% ****** passes..."). Treat both as provisional; do not polish further until
% ****** that rewrite lands.
Responses were anonymized before analysis: names and identifiers were dropped,
and each response carries a term-and-respondent code. From the 873 open-text
responses I derived a coded corpus. The codebook was fixed before coding
began, and its categories come from the dissertation's own analytic
apparatus. It contains roughly forty codes organized into eight families,
each family collecting the codes for one kind of signature: forms of boredom;
flights away from the platform; the texture of reported time; the station of
the actualization chain a response addresses; the involvement channels a
response names; signatures of rhetorical incorporation; student suggestions;
and an open family for codes that emerged during coding. Coders read
each response segment by segment, a segment being the clause or sentence
carrying one codable idea, and recorded a set of codes per response, with
verbatim evidence spans retained for every application in the load-bearing
families. Coding was performed by AI coders working blind under fixed written
instructions; I adjudicated every disputed application myself.\footnote{All
coding passes in the final corpus were performed by Anthropic's Claude Fable 5
model, each pass blind and independent under written coder instructions
reproducing the codebook. Three batches were initially single-coded by a
smaller model (Claude Haiku 4.5) for economy; a same-instrument comparison
measured its agreement with Fable 5 at mean Jaccard .51--.67 against the .878
same-model benchmark, and I discarded those passes and had the batches
recoded, so no Haiku coding remains in the corpus. Per-code agreement on the
double-coded subsample ranged from 50\% to 100\%, weakest on union-merged
mention codes and strongest on the load-bearing families. Every cross-pass
dispute on a load-bearing code reached me for manual ruling, and I
subsequently reviewed the apparatus station by station---every rule, boundary,
and disputed application. A dedicated two-pass recount of the
shared-involvement evidence followed the same protocol (93\% exact agreement).
All pass files, comparison reports, and adjudication records are retained.}
All 873 responses were coded at least once, and a 143-response
subsample---drawn deliberately from the term the codebook had least seen---was
coded twice by independent blind passes, yielding a mean Jaccard agreement of
.878 over full code sets, the set-overlap measure appropriate to multi-label
coding (Krippendorff's $\alpha$ with MASI distance, .82). Merging the two
passes was deterministic and asymmetric: mention-level families merged by
union, the analytically load-bearing families by intersection, with the four
residual disputes resolved by hand; a subsequent line-by-line review of the
whole apparatus added three further manual corrections. Throughout, codes name
signatures in student talk, never certified inner states.

The corpus has four standing limits.
It is a single self-reported channel: no physiological, behavioral, or
performance measure accompanies it, so claims about inner states rest on what
students chose to write. It is retrospective: students answered at the end of
the term, so the corpus preserves remembered experience, not experience in
progress. It supports no learning-outcome claims: the course's content items
sat at ceiling, and nothing here measures whether one format taught better than
another. Term-to-term comparisons carry a composition caveat: the three cohorts
differ in size and possibly in makeup, so cross-term movements are read as
trends, not tests. What the corpus does provide, at a scale the pilot could
not, is 241 students' own accounts of entering---or failing to enter---the same
virtual world.
% --- M2.3 The method of reading — BLOCKED by Open decision #1 (arc placement;
%     deferred until lab-pole O-9 processing). Do not draft before the fork is
%     chosen.
